%! TEX program = pdflatex
% WARNING: this is a generated file.
%
% Please do not edit this file directly. 
% - If you want to update the medatata of the paper (title, authors, abstract), please
%   edit the `paper-meta.yaml` file in the root of the repository.
% - If you want to update the content of the paper, please edit the latex files
%   in the `src` directory.
% - If you want to update the template itself (e.g., change the layout), please
%   edit the `templates/plain-article.tex` file instead.
\documentclass[9pt,a4paper,twosided]{article}

% we setup a custom geometry because the default one is too narrow
\usepackage{geometry}
\geometry{margin=3.5cm}

% utf-8 for old systems
\usepackage[utf8]{inputenc}
\usepackage[T1]{fontenc}

% babel for language settings
\usepackage[english]{babel}

% microtype for better typography
\usepackage{microtype}


% Configure unicode support for xelatex
\ifxetex
  \usepackage{fontspec}
\fi


% math packages
\usepackage{amsmath,amsthm,amssymb,stmaryrd,thmtools,upgreek}

% configure some theorems
\newtheorem{theorem}{Theorem}
\newtheorem{lemma}[theorem]{Lemma}
\newtheorem{corollary}[theorem]{Corollary}
\newtheorem{proposition}[theorem]{Proposition}
\newtheorem{conjecture}[theorem]{Conjecture}
\newtheorem{assumption}[theorem]{Assumption}
\theoremstyle{definition}
\newtheorem{definition}[theorem]{Definition}
\newtheorem{remark}[theorem]{Remark}
\newtheorem{example}[theorem]{Example}
\newtheorem{fact}[theorem]{Fact}
\newtheorem{claim}[theorem]{Claim}


% claim and claim proof environment
\def\cqedsymbol{\ifmmode$\lrcorner$\else{\unskip\nobreak\hfil
		\penalty50\hskip1em\null\nobreak\hfil$\lrcorner$
		\parfillskip=0pt\finalhyphendemerits=0\endgraf}\fi} 
\newcommand{\cqed}{\renewcommand{\qed}{\cqedsymbol}}

\newenvironment{claimproof}[1][Proof.]{\begin{proof}[#1]}{\renewcommand{\qed}{\hfill$\blacksquare$}\end{proof}}
\newcommand{\claimqedhere}{\renewcommand{\qedsymbol}{\hfill$\blacksquare$}\qedhere}



% graphics packages
\usepackage{graphicx}
\usepackage[obeyclassoptions,mode=tex]{standalone}
\usepackage{tikz}
\usetikzlibrary{backgrounds}
\usetikzlibrary{shapes.geometric}
\usetikzlibrary{decorations.markings}
\usetikzlibrary{positioning}
\usetikzlibrary{automata}
\usetikzlibrary{tikzmark}
\usetikzlibrary{patterns}
\usetikzlibrary{arrows}
\usetikzlibrary{calc}

\tikzset{every state/.style={minimum size=1pt}}
\usepackage{tikz-cd}

% ornaments
\usepackage{pgfornament}


% links inside the document
\usepackage{hyperref}
\usepackage[capitalise,noabbrev,nameinlink]{cleveref}
\Crefname{assumption}{Assumption}{Assumptions}
\crefname{assumption}{Assumption}{Assumptions}
\Crefname{conjecture}{Conjecture}{Conjectures}
\crefname{conjecture}{Conjecture}{Conjectures}
\Crefname{fact}{Fact}{Facts}
\crefname{fact}{Fact}{Facts}
\Crefname{remark}{Remark}{Remarks}
\crefname{remark}{Remark}{Remarks}
\Crefname{example}{Example}{Examples}
\crefname{example}{Example}{Examples}
\Crefname{claim}{Claim}{Claims}
\crefname{claim}{Claim}{Claims}

\usepackage[electronic,hyperref,xcolor,cleveref]{knowledge}
\knowledgeconfigure{notion}



% Tables 
\usepackage{booktabs}
\usepackage{varwidth}

% Algorithms
\usepackage{algorithm2e}
\Crefname{algocfline}{Algorithm}{Algorithms}
\crefname{algocfline}{Algorithm}{Algorithms}
\crefname{algocf}{Algorithm}{Algorithms}
\Crefname{algocf}{Algorithm}{Algorithms}

% Packages for macro definitions
\usepackage{xparse}
\usepackage{xpatch}
\usepackage{tokcycle}
\usepackage{ifthen}

% Proofs
\usepackage{bussproofs}

% Colors 
\usepackage{ensps-colorscheme}
% parametrize knowledge colors and style
% intro => A4 + emph 
% kl    => normal color, normal size not emph
\knowledgestyle{intro notion}{color={A2}, emphasize}
\knowledgestyle{notion}{color={B1}}
\hypersetup{
    colorlinks=true,
    breaklinks=true,
    linkcolor=A4,
    citecolor=A4,
    urlcolor=A4,
    filecolor=A4,
}

% Knowledge logo
\newcommand{\klogo}{%
\begin{tikzpicture}[scale=0.2,line/.style={draw, line width=0.2pt, line cap=round, line join=round}]
\coordinate (A00) at (0,0);
\coordinate (A01) at (0,1);
\coordinate (A10) at (1,0);
\coordinate (B10) at (1,0.2);
\coordinate (B01) at (0.2,1);

\coordinate (C01) at (0.4,0.7);
\coordinate (C10) at (0.7,0.4);
\coordinate (C12) at (0.4,1.2);
\coordinate (C21) at (1.2, 0.4);
\coordinate (C22) at (1.2, 1.2);

\coordinate (D00) at (C10);
\coordinate (D01) at (0.8,0.5);
\coordinate (D10) at (0.8,0.3);

\coordinate (E01) at (0.3,0.7);
\coordinate (E10) at (0.5,0.7);

\draw[line] (B01) -- (A01) -- (A00) -- (A10) -- (B10);
\draw[line] (C01) -- (C12) -- (C22) -- (C21) -- (C10);

\draw[line] (D01) -- (D00) -- (D10);
\draw[line] (E01) -- (E10);

\end{tikzpicture}%
}



% we include whatever the user wants to include in the header
\usepackage{ensps-colorscheme}
\usepackage{todonotes}
\newcommand{\mael}[1]{\todo[color=violet!40]{#1}}
\newcommand{\maelin}[1]{\todo[inline,size=\normalsize,color=violet!40,caption={}]{Mael: #1}}

% we include libraries (tex files) usually written in the `lib` directory
\input{lib/aliaume.tex}
\input{lib/maths.tex}
\input{lib/knowledges.kl}

% We include the title and author information based on the 
% `paper-meta.yaml` file.
 
\title{Well-quasi-ordered classes of bounded clique-width}

\author{
Maël Dumas\thanks{University of Warsaw}
 \and
Aliaume Lopez\thanks{INP Bordeaux, LaBRI, CNRS}
}

% For the date, we first check if the user has provided a date,
% and otherwise use the git meta inforamtion (if available).
      \date{\today}
  
\newcommand{\repositoryUrl}{\url{https://github.com/AliaumeL/clique-width-labelled-wqo}}

\knowledge{notion}
 | kl-usage

% Now, we create the document itself.
\begin{document}
% Generate the title page
\maketitle
% Print the abstract
\begin{abstract}
    We study classes of graphs with bounded clique-width that are well-quasi-ordered by the induced subgraph relation, in the presence of labels on the vertices. We prove that, given a finite presentation of a class of graphs, one can decide whether the class is labelled-well-quasi-ordered. This solves an open problem raised by Daligault, Rao and Thomassé in 2010, and answers positively to two conjectures of Pouzet in the restricted case of bounded clique-width classes. Namely, we prove that being labelled-well-quasi-ordered by a set of size 2 or by a well-quasi-ordered infinite set are equivalent conditions, and that in such cases, one can freely assume that the graphs are equipped with a total ordering on their vertices. Finally, we provide a structural characterization of those classes as those that are of bounded clique-width and do not existentially transduce the class of all finite paths.
\end{abstract}

\paragraph{Keywords:}
    well-quasi-ordering, clique-width, automata theory, monoids, factorization forests, gap embedding.
\paragraph{ACM CSS 2012:}
Formal languages and automata theory;
Graph theory.
\paragraph{Repository:} \repositoryUrl

\paragraph{Acknowledgments.}
The authors would like to thank the Zygmunt Zaleski fundation for supporting the authors during the preparation of this work.
Maël Dumas was supported by the ERC project BUKA (n° 101126229) (during his employment in BUKA) and project BOBR (during his employment in BOBR) that received funding from the European Research Council (ERC) under the European Union's Horizon 2020 research and innovation programme (grant agreement No.~948057). 


\bigskip

\klogo\ This document uses \href{https://ctan.org/pkg/knowledge}{knowledge}:
\kl[kl-usage]{notion} points to its \intro[kl-usage]{definition}. 

% Include the content of the paper
\input{src/introduction}
\input{src/prelims}
\input{src/ramseyan}
\input{src/badpaths}
\input{src/maintheorems}
\input{src/interpreting}
\section{Concluding Remarks}
\label{sec:conclusion}

We have proven that for \kl{hereditary classes} of graphs of \kl{bounded
clique-width}, being \kl{$\forall$-well-quasi-ordered} is a very robust notion,
that enjoys many structural characterisations, and where most conjectures from
the literature on labelled well-quasi-orderings hold true. We consider these
results as a first step towards understanding the landscape of labelled
well-quasi-orderings of graphs and other relational structures.
Note that we proved a very strong dichotomy result for hereditary
classes of graphs of bounded clique-width: either they are not
\kl{$2$-well-quasi-ordered} or the \kl{induced subgraph} relation on them is
simpler than the \kl{gap-embedding} relation on trees. We conjecture that this
dichotomy extends to all hereditary classes of graphs: there is nothing beyond
(nested) trees.
We strongly believe that our techniques can be adapted to the case
of finite relational structures over a finite relational signature. Another
interesting direction would be to consider non-hereditary classes of graphs of
bounded clique-width. In this case, we also believe that our techniques can be
adapted.

A fundamental limitation of our work is that we assume classes to be of
\kl{bounded clique-width}. We believe that classes of graphs that are
\kl{$2$-well-quasi-ordered} are necessarily of \kl{bounded clique-width} (see
\cref{cwqo:conj}), and a first step would be to show that every
such class is \emph{monadically dependent}, a notion originating from 
model theory \cite{baldwin85}
that has been recently investigated in finite model theory 
with the belief that it characterises tractable model checking for 
first-order formulas (see for instance \cite{dreier24} and \cite{dreier2024flipbreakability}).
We actually conjecture a stronger statement, that does not involve labelling
the graphs of the class.

\begin{conjecture}
  \label{conj:wqo-mon-dep}
  Every \kl{hereditary class} of graphs that is 
  \kl{well-quasi-ordered} by the \kl{induced subgraph} relation is \emph{monadically dependent}.
\end{conjecture}

Even in the case of classes having \kl{bounded clique-width}, there are 
several unanswered conjectures that are refinement of our results.
For instance, can we
\kl{existentially interpret} all paths in every \kl{hereditary class} of graphs
that is not \kl{$\forall$-well-quasi-ordered} (and of \kl{bounded
clique-width}). An argument similar to \cref{lem:transduction-not-wqo} shows
that if a class \kl{existentially interprets} all paths, then it is not
\kl{$2$-well-quasi-ordered}, but the converse direction remains open.

Similarly, we focused on classes that are \kl{$2$-well-quasi-ordered}
as our simplest possible labelling. However, there are other 
ways to add information to the graphs: one can distinguish $n$ vertices
(adding $n$ constants), or one can assume that the labelling 
itself is ordered by $a \leq b$ ; adapting the notion of embedding in each case. 
It is clear that a class that is 
\kl{$2$-well-quasi-ordered} will be \kl{well-quasi-ordered} when
finitely many constants are added to distinguish nodes 
\cite[Lemma 5.2]{braunfeld21}, and that it will still be 
\kl{well-quasi-ordered} when \kl{freely labelling} the class
with two comparable labels $a \leq b$. We believe that the converse holds.

\begin{conjecture}
  Every \kl{hereditary class} of graphs that is 
  \kl{well-quasi-ordered} by the \kl{induced subgraph} relation 
  with \emph{two} additional constants
  is \kl{$2$-well-quasi-ordered}. 
\end{conjecture}

% Include the bibliography
\bibliographystyle{plainurl}
\bibliography{papers.bib}

% If there are any appendices, we include them here.
\appendix
\clearpage 
\section{Proofs of \cref{sec:ramseyan}}


\begin{proofof}{fact:gap-embedding-gaps}
  This is proven by induction on the length of the path 
  from $u$ to $v$ in $\t_1$. The base case is exactly 
  \cref{item:gap-embedding:edge} of the definition of \kl{gap-embedding}.

  For the inductive case, let $u, v$ be two nodes in $\t_1$
  such that $\spt_1(u:v) > k$ and $\spt_1(v) = k$ for some $k \in \set{1, \dots, N}$.
  Let $w$ be the \kl(tree){parent} of $v$ in $\t_1$.
  Since $\spt_1(u:v) > k$, we have $\spt_1(w) > k$.
  By induction hypothesis, we have $\spt_2(h(u):h(w)) \geq k$.
  Furthermore, by definition of \kl{gap-embedding}, we have $\spt_2(h(w):h(v)) \geq k$.
  From there, we conclude that $\spt_2(h(u):h(v)) \geq k$.
\end{proofof}

\section{Proofs of \cref{sec:bad-patterns}}

\begin{proofof}{lem:finitely-many-boughs}
  Let $\t = C[B]$ be a tree obtained by plugging a \kl{bough} $B$ of level
  $k$ inside a \kl(bough){context} $C[\square]$. Our goal 
  is to extract finitely many parameters from $B$ 
  (independently of $C[\square]$) such that any other \kl{bough} $B'$ of level $k$
  with the same parameters is \kl(bough){compatible} with $B$.
  These parameters will be elements of $M$ or tuples thereof,
  and since $M$ is finite, we will immediately conclude.

  To that end, let us fix a \kl(bough){context} $C[\square]$. Assume that $\t =
  C[B]$ is \kl{forward Ramseyan}: we want to ensure that $\t' = C[B']$ is
  \kl{forward Ramseyan} of any $B'$ with the same parameters. 
  Recall that $\t'$ is \kl{forward Ramseyan} if for
  every branch in $\t'$, for every level $l$, for every \kl{$l$-neighbourhood}
  in this branch, and for every four nodes $x \treelt y$ and $x' \treelt y'$ in
  this \kl{$l$-neighbourhood}, we have $\tlbl{\t'}{x}{y} \cdot
  \tlbl{\t'}{x'}{y'} = \tlbl{\t'}{x}{y}$.

  First, remark that \kl{$l$-neighbourhoods} with $l < k$ inside $\t'$ are fully
  contained either in $B'$, or $\aTree_{\mathsf{root}}$, or
  $\aTree_{\mathsf{left}}$,  or $\aTree_{\mathsf{right}}$. As a consequence,
  nothing needs to be checked on such \kl{$l$-neighbourhoods} since $B'$
  is assumed to be locally \kl{forward Ramseyan}.
  More generally, the only possible failures of the \kl{forward Ramseyan}
  condition in $\t'$ must involve nodes that are not all contained
  in one of these four subtrees.

  Now, let us consider \kl{$l$-neighbourhoods} with $l \geq k$.
  If $l \geq k$, then the only \kl{$l$-neighbourhood} of a branch that
  is not fully contained in one of the four subtrees crosses $B'$. Let us 
  distinguish several cases based on the shape of the branch with respect 
  to the \kl(bough){context} $C[\square]$ and the \kl{bough} $B'$.
  \begin{enumerate}
    \item The branch containing the \kl{$l$-neighbourhood} starts in 
      $\aTree_{\mathsf{root}}$ (possibly in $\square$), goes through $B'$, and ends in 
      $\aTree_{\mathsf{left}}$ or $\aTree_{\mathsf{right}}$.
    \item The branch containing the \kl{$l$-neighbourhood} starts in 
      $\aTree_{\mathsf{root}}$ (possibly in $\square$), goes through $B'$, and ends in $B'$.
  \end{enumerate}

  In the first case, consider two nodes $x \treelt y$ and $x' \treelt y'$ in
  the \kl{$l$-neighbourhood}. Remark that the value $\tlbl{\t'}{x}{y}$ can be
  decomposed into a part before entering $B'$, a part inside $B'$, and a part
  after leaving $B'$. By construction, the part inside $B'$ is one of the
  idempotent elements labelling the \kl(bough){backbone} of $B'$. Hence, to
  ensure that $\t'$ is \kl{forward Ramseyan} in this case, it suffices to check
  finitely many equations involving those idempotent elements (note that the
  number of equations depends on the shape of $C[\square]$, but not 
  the number of elements of $M$ extracted from $B'$).

  In the second case, consider two nodes $x \treelt y$ and $x' \treelt y'$ in
  the \kl{$l$-neighbourhood}. Here, the value $\tlbl{\t'}{x}{y}$ can be
  decomposed into a part before entering $B'$, a part from $b_0'$ to the first
  node $z$ in $B'$ of level $l$ (which is not on the \kl(bough){backbone} of
  $B'$ if $l > k$), and a part from $z$ to $y$. Similarly, the equations
  ensuring that $\t'$ is \kl{forward Ramseyan} in this case can be expressed
  using the values $\tlbl{\t'}{b_0'}{z}$ for every node $z$ in $B'$, and the
  values $\tlbl{\t'}{z}{y}$ for every node $y$ in $B'$ in the same
  \kl{$l$-neighbourhood} as $z$ (for some branch). Note that one only cares
  about the existence of pairs of values $(\tlbl{\t'}{b_0'}{z},
  \tlbl{\t'}{z}{y})$ for nodes $z$ and $y$ in the same \kl{$l$-neighbourhood},
  and there are only finitely many such pairs.

  We have proven that, assuming that $\t = C[B]$ is
  \kl{forward Ramseyan}, it suffices to check finitely many parameters of $B'$
  to ensure that $C[B]$ is \kl{forward Ramseyan}. Since $M$ is finite, there
  are finitely many possible values for these parameters, which concludes the
  proof.
\end{proofof}

\section{Proofs of \cref{sec:hereditary-classes}}

\begin{proofof}{lem:recognizing-perfect-boughs}
  The proof follows the same pattern as 
  \cite[Lemma 19]{LOPEZ24}, but adapts it from words to trees. 
  Assume that $B$ is a \kl{perfect bough} of level $k$ inside of a tree
  $\t = C[B]$.
  This gives us a compatible \kl{bough} $H$ of level $k$
  and an embedding $h \colon \someInterp(C[B]) \to \someInterp(C[H])$
  satisfying the conditions of \cref{def:perfect-bough}.

  Let us consider five copies of $B$, called $B_\mathsf{left}$,
  $B_\mathsf{mid-1}$, $B_\mathsf{mid-2}$, $B_\mathsf{mid-3}$, and
  $B_\mathsf{right}$, and let us define $B^5$ to be the \kl{bough} obtained by
  concatenating these five copies merging $b_n$ of $B_\mathsf{left}$ with $b_0$
  of $B_\mathsf{mid-1}$, and merging $b_n$ of $B_\mathsf{mid-1}$ with $b_0$ of
  $B_\mathsf{mid-2}$ and so on. We claim that $B$ and $B^5$ are \kl{compatible
  boughs}. 

  Let us define a mapping $k \colon \someInterp(C[B]) \to
  \someInterp(C[B^5])$
  as follows: leaves $x$ of $B$ are mapped to their
  corresponding leaf in $B_\mathsf{left}$ or $B_\mathsf{right}$ depending on
  whether $h(x)$ is to the left or right of the untouched \kl{bough blocks} of
  $H$. For leaves $x$ of $C[\square]$, we let $k(x) = x$. 
  Thanks to \cref{fact:bough-replacement}, we know that
  $k$ preserves edges between leaves of $C[\square]$.
  Furthermore, because $h$ is an embedding that preserves the 
  \kl{bough type} of leaves in $B$, and because \kl{bough types} are enough
  to compute the presence of an edge between a leaf of $B$ and a leaf of
  $C[\square]$ (or between two leaves of $B$), we deduce that $k$ is also an
  embedding.

  What we have shown is that if $B$ is a \kl{perfect bough}, then there exists
  a mapping $k'$ from leaves of $B$ to $\set{ \mathsf{left}, \mathsf{right} }$
  such that there is an edge between two leaves $x$ and $y$ of $B$ in the graph
  $\someInterp(C[B])$ if and only if the \kl{bough types} of $x$ and $y$ are
  such that there is an edge between $x_{k'(x)}$ and $y_{k'(y)}$ in the graph
  $\someInterp(C[B^5])$. One can guess such a mapping $k'$ in Monadic Second-Order
  logic, and then verify the required conditions on edges between leaves of $B$.
  As a consequence, one can write an $\MSO$ formula that holds if and only if
  $B$ is a \kl{perfect bough}.
\end{proofof}

\begin{proofof}{lem:hereditary-perfect-boughs}
  The proof follows the same pattern 
  as \cref{lem:finitely-many-boughs,lem:finitely-many-contexts},
  and continues on the ideas of \cref{lem:recognizing-perfect-boughs}.
  The idea is that if a bough $B$ is \kl(bough){perfect}, then 
  one can witness it by using the \kl{compatible bough} $B^5$ 
  defined by gluing five copies of $B$ together, because we noticed
  that perfectness essentially depends on the \kl{bough types} of the leaves
  of $B$. 
  Another consequence of this analysis is that if one defines an equivalence relation between
  \kl{bough blocks} of a given \kl{bough} $B$ based on the set of 
  available \kl{bough types} of its leaves, then 
  one can always replace a \kl{bough block} by another one in the same
  equivalence class without changing the \kl(bough){perfectness} of $B$.
  We conclude because this equivalence relation has finitely many classes.
\end{proofof}



\end{document}
